%======================================================================
\section{Invariant decomposition of the constant-profile \texorpdfstring{\(hh\)}{hh} poles (new in v6)}
\label{sec:structures}
%======================================================================

On the axisymmetric lattice, write \(h\) for the longitudinal
\(2\times2\) block, \(\operatorname{tr}h=\eta^{ab}h_{ab}\),
\(h\cdot h'=h^{ab}h'_{ab}\), and
\(php=p^ap^bh_{ab}\) for \(a,b\in\{0,3\}\).
For two-dimensional symmetric matrices the Cayley--Hamilton
identity removes \(p\,hh'\,p\).
A symmetric \(SO(1,1)\)-invariant bilinear that is quadratic in
\(p\) is therefore a combination of six structures:
\begin{align*}
S_1&=P^2\operatorname{tr}h\operatorname{tr}h', &
S_2&=P^2\,h\cdot h', &
S_3&=(php)\operatorname{tr}h'+\operatorname{tr}h\,(ph'p),\\
S_6&=P^2(f_0\operatorname{tr}h'+\operatorname{tr}h\,f_0'), &
S_7&=(php)f_0'+f_0(ph'p), &
S_8&=P^2f_0f_0'.
\end{align*}

\begin{proposition}[Structure coefficients]
\label{prop:structures}
\textsc{[exec]}
Every executed constant-profile \(hh\) pole decomposes uniquely as
\(\sum_ib_iS_i\).
In units of \(\ka^2/\pi\), the coefficients are the following
branch-uniform linear forms in \((c_0,c_1,c_2)\):
\begin{center}
\small
\begin{tabular}{@{}lcccccc@{}}
\toprule
Graph & \(b_1\) & \(b_2\) & \(b_3\) & \(b_6\) & \(b_7\) & \(b_8\)\\
\midrule
D01 & \(-\tfrac1{16}c_0\) & \(\tfrac18c_0\) & \(0\) &
\(-\tfrac18c_0-\tfrac14c_1-\tfrac12c_2\) &
\(\tfrac14c_0+\tfrac12c_2\) &
\(\tfrac14c_0-\tfrac12c_1+\tfrac14c_2\)\\
D02 & \(0\) & \(0\) & \(0\) & \(W\) & \(-W\) & \(-W\)\\
D03 & \(\tfrac1{16}c_0\) & \(-\tfrac18c_0\) & \(0\) &
\(\tfrac38c_0\) & \(-\tfrac14c_0\) &
\(-\tfrac14c_0-\tfrac12c_1\)\\
D11 & \(0\) & \(0\) & \(0\) & \(0\) & \(0\) & \(\tfrac18c_0\)\\
D12 & \(0\) & \(0\) & \(0\) & \(0\) & \(0\) & \(0\)\\
\midrule
Sum & \(0\) & \(0\) & \(0\) &
\(\tfrac34(c_0+c_1)\) & \(-\tfrac12c_0-c_1\) &
\(-\tfrac38c_0-2c_1-\tfrac14c_2\)\\
\bottomrule
\end{tabular}
\end{center}
Here \(W=\tfrac12(c_0+2c_1+c_2)\).
The CR fit determines the forms uniquely.
On F, where \(c_2^\Fr=7c_0+6c_1\) makes the three coefficients
linearly dependent, the same forms reproduce every entry
(\loc{logs/v6/hh\_structures.json}).
UNIT (\(c\equiv1\)) gives the sum
\((0,0,0,\tfrac32,-\tfrac32,-\tfrac{21}8)\).
\end{proposition}

\begin{corollary}[Gauge and branch content of the complete \(hh\) sum]
\label{cor:hhcontent}
\leavevmode
\begin{enumerate}[label=\textup{(\roman*)}]
\item The purely longitudinal structures cancel:
\(b_1=b_2=b_3=0\).
This cancellation is between D01 and D03.
\item The only gauge-variant combination on the lattice is
\[
b_6+b_7=\tfrac14(c_0-c_1)\frac{\ka^2}{\pi}
=-\frac{\ka^2V}{4\pi}
=-\frac{\ka^2(\alpha-1)}{8\pi},
\]
which is branch independent.
This is exactly the term balanced by the antifield row in
Theorem~\ref{thm:STclosure}.
\item The CR--F difference of the complete \(hh\) sum is
\[
\Delta\widehat B^{hh}=-\frac{\nu\ka^2}{4\pi}P^2f_0f_0'.
\]
It is annihilated by the longitudinal gauge vector
\eqref{eq:gaugevector}, since \(k_T=0\).
The D01 difference of Proposition~\ref{prop:branchtransverse} has
\(S_6,S_7\) parts. They cancel against the D02 difference
\(\tfrac{\nu}{2}(S_6-S_7-S_8)\ka^2/\pi\).
This settles the v5 obligation that D02/D12 must be included in
the branch-transversality check.
\end{enumerate}
\end{corollary}

\begin{remark}[Relation to the 18-dimensional engine map]
Axisymmetric constant-profile data land in the six-dimensional
space above.
The programme's complete \(hh\) map is 18-dimensional
\cite{CheungConeV4,CheungConeV5}; it is computed by the
XVI-P engine in its own basis.
That basis also includes transverse-derivative and non-axisymmetric
structures, which constant profiles cannot see
(Proposition~\ref{prop:noninjective}).
The convention map that sends the six coefficients above to their
images in the engine's 18 coordinates is \(\OPEN\).
So, therefore, are the components of the 18-dimensional vector
outside the image.
They are not set to zero.
\end{remark}
